\begin{textsample}{POS Dim 1 – go – Score 41.00 – go/go\_inf\_23.txt}  \label{ex:f1_pos_105}
3.5.
Espaços, Tempos, Quantidades, Relações e Transformações A criança, desde bebê, demonstra curiosidade, interesse em compreender o funcionamento do mundo.
De início percebe e interage com pessoas, objetos e materiais pelas sensações, olhando, pegando, levando a boca, passando no corpo, testando se é duro ou mole, se faz barulho, se deixar cair o objeto qual será a reação dos outros.
Enfim, aprende e se desenvolve em todas as suas dimensões, expressivo-motora, afetiva, linguística, ética, estética e sociocultural ( BRASIL, 2009 ), significando e produzindo sentido às suas ações na relação com os outros.
À medida \textbf{que} desenvolve a fala iniciam -se os questionamentos – “ O \textbf{que} é isso? ”, “ Para \textbf{que} serve? ”, “ É de plástico? ”, “ É de madeira? ”.
Indaga ainda sobre fenômenos naturais e aspectos socioculturais – “ Por \textbf{que} eu tenho \textbf{que} acordar no escuro ( ao se referir ao horário de verão )? ”, “ Como \textbf{que} compra dinheiro? ”, “ Quais eram as brincadeiras de antigamente? ”, “ As abelhas dormem? ”, “ Quantos anos você tem? ”, “ Você já é velhinha? ”, “ Eu também vou morrer? ”, “ Quando eu morrer vou comer formigas? ”.
Essa necessidade de compreender os porquês da vida, \textbf{que} mundo é esse e como ele se organiza faz parte da natureza humana. \textbf{Dependendo} da maneira como os professores escutam as perguntas da criança e as problematizam, fazendo novas indagações e planejando situações de aprendizagens \textbf{que} a desafie a construir explicações para os seus porquês, a criança poderá ter aguçadas, a curiosidade, a criatividade e a criticidade, desenvolvendo \textbf{uma} forma própria de pensar e de se posicionar diante de fatos, situações e conhecimentos.
A BNCC ( BRASIL, 2017 ) apresenta este campo de experiências abordando conceitos \textbf{centrais} \textbf{que} possibilitam compreender como a criança percebe e significa o mundo físico e sociocultural, envolvendo os espaços, os tempos e os conhecimentos matemáticos, conforme segue abaixo : Esse campo tem por finalidade, de forma interligada aos demais campos, possibilitar a construção de noções espaciais, temporais e quantitativas ; o estabelecimento de relações de causa e efeito e a \textbf{percepção} de regularidades, de irregularidades e de transformações no mundo físico e sociocultural.
Para tanto, os processos de apropriação e produção desses conhecimentos de diferentes naturezas devem ser trabalhados, de forma regular e sistemática, a partir de processos investigativos \textbf{que} leve a criança a : observar ; explorar ; estabelecer relações entre fatos, objetos, acontecimentos, noções, conceitos ; levantar hipóteses ; questionar, buscar soluções ; verificar resultados ; \textbf{sistematizar} informações e dados coletados em várias fontes ; construir explicações e registrar por meio de diferentes linguagens o percurso vivenciado.
O Mundo Físico e o Mundo Sociocultural A criança, desde sua concepção, é inserida, participa e \textbf{vive} num espaço específico, o planeta Terra.
Este possui características e fenômenos próprios, como : o nascer e o pôr do sol, as tempestades, os processos de decomposição e germinação, a erosão, entre outros.
Há ainda recursos naturais como o solo, a água, o oxigênio, a energia oriunda do sol, as florestas, os minerais, \textbf{que} permitem a vida e o desenvolvimento de \textbf{inúmeras} espécies de animais, de plantas, de bactérias, de vírus e do próprio \textbf{homem}.
Denomina -se mundo físico esse todo complexo, compreendido a partir de leis \textbf{que} regem seu funcionamento.
O \textbf{homem} como \textbf{uma} das espécies \textbf{que} habita esse mundo na relação com outros \textbf{homens} – no decorrer da história, por diferentes motivos, sobrevivência, formação de grupos, descoberta de novas terras, garantia de conforto e bem-estar, dominação de outros povos – produziu e produz, artefatos, costumes, hábitos, regras, conhecimentos, atitudes, valores, línguas, tecnologias, relações espaciais, noções de tempo, crenças etc., para organizar e estruturar sua vida na terra.
Essas produções de diferentes povos, grupos sociais e culturas, resultado da invenção e do trabalho humano, constituem o legado deixado de \textbf{uma} geração para outra, como tradição ou patrimônio da \textbf{humanidade}.
Todos esses conhecimentos, relações, saberes constituem o mundo sociocultural.
Assim, é na interdependência e na \textbf{influência} mútua entre o mundo físico e o mundo sociocultural \textbf{que} a vida acontece, permitindo a cada \textbf{sujeito} \textbf{que} chega a esse planeta se apropriar do \textbf{que} já foi produzido e, ao mesmo tempo, deixar sua \textbf{marca} e introduzir novos elementos.
Os conceitos \textbf{centrais} relacionados aos mundos, físico e sociocultural, \textbf{que} perpassam por este campo de experiências e \textbf{que} serão destacados neste texto, são : espaços, tempos e conhecimentos matemáticos.
Espaços A criança está inserida em variados espaços, desde \textbf{uma} escala local - a rua, o bairro, o município - até \textbf{uma} escala global - o estado, o país, o continente, o planeta Terra.
Auxiliá -la a se localizar e a ler esses diversos espaços possibilita \textbf{uma} maior compreensão e atuação dela no mundo.
Essa é \textbf{uma} das funções das instituições de Educação Infantil.
Os espaços são entendidos como as diferentes paisagens e lugares existentes e incluem tanto os naturais como os construídos ou modificados pelos \textbf{homens} ao longo do tempo.
Os espaços se constituem a partir das relações estabelecidas entre os \textbf{sujeitos} no e com o mundo.
Para Horn ( 2017, p. 18 ) > [... ] a construção do espaço é eminentemente social e se entrelaça com o tempo de forma indissolúvel, congregando, de forma simultânea, diferentes \textbf{influências} mediatas e \textbf{imediatas} advindas da cultura e do meio em \textbf{que} seus \textbf{atores} estão inseridos.
As \textbf{influências} mediatas e \textbf{imediatas} realizadas pelo \textbf{homem} no espaço por meio do trabalho têm como suporte as tecnologias, entendidas como técnicas e procedimentos utilizados para produção de variadas ferramentas, objetos e instrumentos \textbf{que} permitem o ser humano intervir no mundo, alterando às vezes, inclusive, a sua ordem natural.
Constituem -se em tecnologias os diferentes objetos \textbf{que} fazem parte do nosso cotidiano, talheres, móveis, vestimentas etc. e os \textbf{que} já são comumente associados a elas, como : computadores, tablets, smartphones, entre outros.
Elas estão presentes desde os primórdios da \textbf{humanidade}, modificando ao longo do tempo as formas, de \textbf{viver} e de ser, dos \textbf{sujeitos}.
Atualmente, é impossível \textbf{viver} sem tecnologia, principalmente, as digitais²¹.
Elas têm revolucionado o mundo do trabalho, do entretenimento, da saúde etc. e das relações interpessoais \textbf{que}, em determinadas situações, tornam -se virtuais.
Portanto, não é possível ignorar sua \textbf{influência} no cotidiano e a necessidade de se ter \textbf{uma} visão crítica sobre o seu uso, \textbf{uma} vez \textbf{que} podem servir tanto para entreter, garantir conforto e qualidade de vida, como para construir armas químicas, incentivar o consumo exacerbado impulsionado pelas inovações, dominar a natureza e seus fenômenos, degradar o meio ambiente, entre outros.
Essa realidade suscita \textbf{uma} nova postura da sociedade, das instituições educacionais e dos seus profissionais, no sentido de criarem contextos \textbf{que} possibilitem o desenvolvimento de \textbf{um} pensamento crítico e analítico a respeito das diferentes tecnologias, reconhecendo as suas possibilidades e limites para a manutenção da vida das mais diversas espécies na Terra. ²¹ Segundo Coscarelli e Ribeiro ( 2011 ) a palavra digital quando se refere à computação tem \textbf{um} sentido amplo está relacionada a \textbf{um} modo de processar, transferir ou guardar informações.
Geralmente, os \textbf{homens} se relacionam com a natureza numa perspectiva utilitarista e de superioridade em relação aos demais seres vivos, não tendo consciência das \textbf{inúmeras} interferências \textbf{que} têm provocado no mundo físico.
Por isso, é \textbf{urgente} a necessidade do ser humano se perceber como parte integrante da natureza, como mais \textbf{uma} espécie entre tantas outras \textbf{que} \textbf{vivem} nesse Planeta.
De acordo com Elias ( 1998, p. 12 ) os \textbf{homens} deveriam compreender \textbf{que} “ [... ] não \textbf{vivem} isolados, mas se inserem no devir da natureza, e \textbf{que}, em função de sua natureza específica, compete a eles a responsabilidade de se encarregarem, em seu próprio interesse, dessa relação ”.
Portanto, cabe ao \textbf{homem} respeitar, proteger, conservar e preservar a natureza para \textbf{que} possa garantir o equilíbrio do meio ambiente e da sua vida na Terra.
Só é possível respeitar, proteger, conservar e preservar aquilo \textbf{que} se conhece e \textbf{que} se considera importante.
Assim, é fundamental permitir e possibilitar à criança estabelecer \textbf{uma} relação próxima da natureza, conhecendo suas características, fenômenos e recursos.
No Estado de Goiás, os tempos de chuva e de seca, o cerrado com sua fauna e flora tão diversa.
Também, é necessário desenvolver atitudes mais sustentáveis em ações cotidianas, por exemplo : recusar o canudo para beber algum líquido, fechar a torneira ao escovar os dentes, apagar as luzes quando sair de \textbf{um} espaço, trocar brinquedos no lugar de comprar novos, entre outras ações \textbf{que} parecem pequenas, mas \textbf{que} podem fazer a diferença.
Nesse processo de (re)conexão do \textbf{homem} com a natureza, Tiriba ( 2010 ) \textbf{pontua} \textbf{que} é necessário olhar com mais atenção a forma como os povos indígenas, os agricultores familiares, os extrativistas, os pescadores artesanais, os \textbf{ribeirinhos}, os assentados e acampados da reforma agrária, os quilombolas, os caiçaras, os povos da floresta se relacionam com o mundo físico, para \textbf{que} se possa aprender com eles outras possibilidades de conviver com a natureza, numa perspectiva de respeito e de \textbf{percepção} dela como \textbf{um} organismo vivo e interconectado.
Porém, as questões ambientais não são fáceis de solucionar.
Ao contrário, são complexas e envolvem aspectos legais, políticos, econômicos, éticos, culturais e científicos de forma interligada numa escala local e global.
Por isso, a necessidade de Convenções, Acordos, Tratados e Legislações \textbf{que} regulamentem usos, posturas e atitudes diante do meio ambiente.
No \textbf{que} se refere ao âmbito educacional, no Brasil, a Lei 9.795 de 1999 do Ministério da Educação e do Meio Ambiente institui a Política Nacional de Educação Ambiental e compreende a educação ambiental, como > os processos por meio dos quais o indivíduo e a coletividade constroem valores sociais, conhecimentos, habilidades, atitudes e competências voltadas para a conservação do meio ambiente, bem de uso comum do povo, essencial à sadia qualidade de vida e sua sustentabilidade. ( BRASIL, 1999, Artigo 1º ) Nessa perspectiva, as instituições de Educação Infantil têm \textbf{um} papel fundamental para o desenvolvimento e o crescimento das crianças mais conectadas à natureza e aos diferentes espaços \textbf{que} convive, bem como, na socialização e produção de conhecimentos científicos e tecnológicos \textbf{que} as subsidiem, assim como às suas famílias e à comunidade onde está localizada a atuarem na prevenção, na identificação e na solução de problemas ambientais.
Como as crianças aprendem a se localizar e a diferenciar os variados espaços em \textbf{que} \textbf{vivem}, assim como os \textbf{que} são \textbf{vistos} apenas em desenhos, vídeos, fotografias etc.?
A criança se apropria do espaço, por meio da \textbf{percepção} de si mesma, \textbf{que} aos poucos se diferencia dos outros, perpassando pela compreensão dela no mundo, até chegar ao entendimento de \textbf{um} espaço \textbf{que} pode ser representado.
Para Smole, Diniz e Cândido ( 2003, p.16 ), “ [... ] a \textbf{percepção} do espaço avança em \textbf{uma} \textbf{direção} marcada por três etapas essenciais : a do \textbf{vivido}, a do percebido e a do concebido ”.
O \textbf{vivido} consiste na ação da criança no espaço, movimentando -se, deslocando -se, organizando e explorando suas dimensões, texturas, temperaturas e objetos com características e propriedades específicas.
O percebido se refere à capacidade de \textbf{lembrar} -se do espaço sem vivenciá -lo fisicamente novamente e o concebido está relacionado com o processo de representação do espaço, em \textbf{que} são estabelecidas relações espaciais a partir de mapas, maquetes, plantas, figuras, coordenadas etc.
Assim, a criança, desde o nascimento, se apropria do espaço por meio da exploração, \textbf{um} dos seus direitos de aprendizagens e desenvolvimento, construindo noções espaciais, por ocupar, mover -se, interagir e \textbf{viver} num espaço.
Dessa forma, quando ela chega às instituições de Educação Infantil já possui diversos saberes e conhecimentos \textbf{que} devem ser ampliados e diversificados por meio de contextos de aprendizagens significativas, como : - deslocamentos pela instituição e ao seu redor para identificar os diferentes espaços suas funções e objetos \textbf{que} os compõem ; - relatos e desenhos de trajetos realizados ; - construção de mapas \textbf{que} permitam a localização de determinados objetos e lugares ; - manuseio, leitura e interpretação de mapas, plantas baixas, globos etc. ; - brincadeiras como caça ao tesouro, encontrar objetos a partir de coordenadas etc. ; - atividades ou brincadeiras nas áreas, interna e externa, \textbf{que} permitam delimitar e organizar o espaço e os objetos de forma autônoma, estabelecendo relações ; - situações do cotidiano \textbf{que} possibilitem identificar as características e as propriedades de objetos, assim como suas funções, estabelecendo relações e comparações entre eles.
Auxiliar a criança a iniciar o desenvolvimento da \textbf{percepção} espacial no sentido de \textbf{conseguirem} ler, se orientar, localizar e representar \textbf{um} espaço é necessário para \textbf{que} elas possam ter \textbf{uma} visão crítica e analítica dos espaços naturais e construídos existentes no mundo, assim como, possuír maior autonomia para sua locomoção, planejamento e organização de diferentes espaços.
Tempos Segundo Elias ( 1998 ), o tempo como é concebido, na atualidade, como \textbf{um} fluxo regular, uniforme e contínuo, dividido em anos, meses, dias, horas, minutos e segundos, medido com instrumentos \textbf{precisos} como o calendário e o relógio ; é \textbf{uma} construção social \textbf{que} surgiu da necessidade de orientar, regular e ordenar a vida e o dia a dia das pessoas.
Em outras épocas ou mesmo em alguns grupos sociais nos dias de \textbf{hoje}, o tempo é compreendido e calculado de forma diferente, tendo como principal referência a natureza e seus fenômenos, como : as estações do ano, as fases da lua, o ritmo das marés, o nascer e o pôr do sol.
Essas referências permitem estabelecer regularidades e organizar ações do cotidiano, como plantio, colheitas, festas e rituais.
Os calendários e os relógios, instrumentos simbólicos \textbf{que} representam a noção do \textbf{que} é o tempo, \textbf{desempenham} a mesma função \textbf{que} os fenômenos naturais, qual seja, organizar a vida coletiva.
Isso é possível porque tanto os objetos criados pelo \textbf{homem}, como a natureza, permitem \textbf{que} as pessoas cumpram determinadas \textbf{tarefas} de forma interligada entre os diferentes membros de \textbf{um} grupo social, como, no caso dos relógios, \textbf{que} propiciam aos sujeitos chegarem com pontualidade no trabalho ou nas instituições educacionais, do calendário, \textbf{que} delimita período de férias e dos fenômenos naturais \textbf{que} estabelecem as épocas de plantio e de colheita, a lua boa para pesca ou corte de cabelo etc.
Dessa forma, os calendários e os relógios foram produzidos em estreita relação e interdependência com o mundo físico.
Como \textbf{pontua} Elias ( 1998, p. 13 ) > Os problemas \textbf{que} os \textbf{homens} procuram resolver, ao medirem a “ duração ”, \textbf{remetem} ao fato de \textbf{que} os grupos humanos estão situados no interior de \textbf{um} conjunto mais vasto do \textbf{que} o formado por eles : o universo natural.
Em toda parte onde se \textbf{opera} com o “ tempo ”, os \textbf{homens} são \textbf{implicados} juntamente com seu meio ambiente, ou seja, com processos físicos e sociais.
Nesse sentido, os instrumentos criados pelos \textbf{homens} surgem de intensa observação e conhecimento de fenômenos da natureza, bem como da aprendizagem e da experiência acumulada pelos \textbf{homens} ao longo da história da \textbf{humanidade}.
Foi necessário o conhecimento elaborado por várias gerações para se chegar numa síntese e generalização do \textbf{que} é o conceito de tempo nos dias de \textbf{hoje}.
Os calendários e os relógios \textbf{que} se equivalem aos movimentos da Terra permitem aos \textbf{sujeitos} situarem a sua vida, no tempo e no espaço, e a terem como referência para estabelecer relações com acontecimentos ligados a outras fases ou períodos da sociedade.
Essa \textbf{correlação}, do tempo com a vida do \textbf{sujeito}, possibilita o estabelecimento de relações entre passado, presente e futuro, a compreensão de fatos e processos sociais específicos - escravidão, revolução industrial etc.- assim como, o entendimento das várias formas de organização dos grupos sociais existentes em outras épocas e na atualidade.
As relações temporais de continuidade e rupturas, permanências e transformações, sucessão e simultaneidade, também auxiliam na compreensão do \textbf{que} é o tempo e se fazem presentes tanto no cotidiano dos \textbf{sujeitos} quanto ao longo da história da \textbf{humanidade}.
Segundo Elias ( 1998 ), em algumas situações os \textbf{homens} esquecem \textbf{que} os calendários e os relógios consistem numa produção humana e naturalizam a noção de tempo, acreditando \textbf{que} sempre foi concebido dessa forma.
Compreender esse percurso histórico e as diferentes formas de entender o tempo, possibilita aos \textbf{sujeitos} estabelecerem \textbf{uma} relação mais reflexiva e consciente na organização do seu tempo, a fim de problematizar \textbf{uma} expressão \textbf{que} comumente se ouve “ o dia deveria ter mais horas, porque o tempo é insuficiente para tantas \textbf{tarefas} a cumprir ”.
De \textbf{que} forma o tempo é naturalizado nas instituições de Educação Infantil?
Quando o tempo de \textbf{uma} instituição em tempo integral, em \textbf{que} as crianças passam de 8 a 10 horas por dia, torna -se insuficiente?
O tempo é naturalizado, nas instituições de Educação Infantil, quando a rotina diária é organizada na mesma lógica e perspectiva do capitalismo, a da produção, em \textbf{que} se acredita \textbf{que} quanto mais atividades forem realizadas pela e com a criança melhor será o seu aprendizado e desenvolvimento.
Dessa forma, as ações pedagógicas são desenvolvidas de forma cronometrada e rápida, porque a criança tem muitas atividades a serem realizadas ao longo do dia.
Nessa perspectiva, não se considera a qualidade das ações e as variáveis \textbf{que} podem interferir no cotidiano de \textbf{uma} instituição, por exemplo : o convite para \textbf{um} passeio ; as necessidades do grupo de crianças naquele dia, quer seja de ficarem mais quietas ou de realizarem movimentos amplos ; o clima chuvoso, muito quente ou frio ; a visita de alguém ; dentre tantas outras, submetendo a criança a \textbf{um} ritmo da vida adulta e da lógica do capital.
A criança não nasce com essa necessidade \textbf{que} se tem, na \textbf{contemporaneidade}, de cronometrar todas as \textbf{tarefas} \textbf{que} são realizadas.
Esse processo de autorregulação, de compreensão e de medição do tempo são aprendizagens complexas \textbf{que} os \textbf{sujeitos} constroem e se apropriam a partir das vivências \textbf{que} lhes são oportunizadas, das relações e das concepções \textbf{que} se tem sobre o tempo, no contexto do qual fazem parte.
Por isso, o desenvolvimento de noções temporais pela criança está diretamente relacionado à maneira como as pessoas com quem ela convive concebem o tempo.
Em contraposição à lógica apresentada de naturalização do tempo, pode -se concebê -lo como \textbf{um} elemento \textbf{que} organiza o cotidiano e \textbf{que} propicia à criança compreender a lógica do dia, numa perspectiva contínua e flexível, em \textbf{que} o inesperado e o inusitado, caso aconteça, é levado em consideração no (re)planejamento do tempo e das ações, assim como o respeito aos diferentes ritmos e as solicitações apresentadas pelo grupo de crianças a partir de diferentes linguagens.
Essa previsibilidade das ações \textbf{que} serão realizadas ao longo do dia num continuum possibilita à criança ter segurança do \textbf{que} está por vir, identificar qual é o seu ritmo e qual é o do da turma, diferenciar momentos do dia e desenvolver a noção de tempo, sem ser somente numa perspectiva cronológica.
O conceito de tempo envolve três dimensões distintas, porém interligadas : - a física \textbf{que} \textbf{diz} respeito às mudanças climáticas, como, quente, frio, úmido, seco, ensolarado, chuvoso ; - a cronológica relacionada à ideia de ano, mês, dia, noite, ontem, \textbf{hoje}, amanhã, horas, minutos e segundos e ; - a subjetiva \textbf{que} se refere à forma como cada \textbf{sujeito} vivencia e incorpora a noção de tempo no seu cotidiano.
Nas instituições de Educação Infantil considerar como cada criança vivencia o tempo e se organiza diante das ações propostas é fundamental para o(a) professor(a) elaborar o seu planejamento.
Para tanto, observar em quais momentos a criança fica mais concentrada no \textbf{que} está fazendo ( parquinho, modelagem com massinha, exploração com tinta guache, brincadeiras na areia, na terra, com pedras e gravetos, experimentos, investigações, brincadeiras de faz de conta ) e em quais atividades ela solicita ao professor mais tempo, lhe possibilita identificar o ritmo da turma, suas preferências e necessidades.
Ações pedagógicas \textbf{que} envolvem o uso e a compreensão do calendário, a construção com as crianças da rotina de forma escrita e com imagens, a aquisição de relógios digitais e analógicos para cada sala e sua utilização no dia a dia, são ações \textbf{que} possibilitam às crianças a se localizarem diante do tempo e a desenvolverem noção temporal.
Conhecimentos Matemáticos Nos diferentes lugares \textbf{que} a criança frequenta e nas práticas sociais das quais participa, os conhecimentos matemáticos \textbf{permeiam} suas ações.
Ao nascer, é medida e pesada para acompanhamento do crescimento com registros no cartão de vacina.
Quando começa a movimentar as mãos e os dedos é incentivada a mostrar com \textbf{um} gesto a idade.
Pessoas perguntam ao responsável se é o primeiro ou o segundo filho.
Os familiares, quando torcem por \textbf{um} time de futebol, adquirem para a criança \textbf{uma} camisa com o número do jogador preferido.
Os professores em ações na área externa contam as crianças para verificarem se estão todas próximas.
Fazem cálculos na frente delas para distribuir porções com a mesma quantidade de alimentos.
Verificam distâncias para realizar passeios no entorno, utilizando os termos, perto, longe, próximo, distante.
Essas vivências são o ponto de partida para ampliação e diversificação dos conhecimentos da criança a partir da articulação com os \textbf{que} já foram produzidos e \textbf{sistematizados} ao longo do tempo pela \textbf{humanidade}.
O trabalho com os conhecimentos matemáticos, a partir da compreensão da criança como \textbf{sujeito} ativo, propositivo, capaz e \textbf{que} tem direito de aprender participando, explorando, convivendo, conhecendo -se, brincando e se expressando, \textbf{implica} em planejar situações de aprendizagens significativas \textbf{que} envolvam contextos do mundo real em \textbf{que} ela : resolva problemas ; aprenda a ler e escrever números ; faça contagens ; brinque com cartas, jogos de tabuleiro ; observe espaços ; localize objetos ; compare medidas ; colete dados ; produza tabelas e gráficos.
Para tanto, de acordo com Monteiro ( 2010 ), as crianças, por meio da resolução de problemas, precisam ter oportunidade para elaborar hipóteses, pesquisar, analisar, inferir, deduzir, utilizar material concreto, realizar cálculo mental e representar suas ações e pensamentos por diferentes tipos de registros, quer seja no coletivo, em pequenos grupos, em duplas ou individualmente.
O desafio das instituições de Educação Infantil é romper com \textbf{uma} perspectiva preparatória para o Ensino Fundamental \textbf{baseada} em ações mecânicas e repetitivas.
Por exemplo, acreditar \textbf{que} a criança desenvolve noções como : em cima, embaixo, entre e ao \textbf{lado}, pintando a localização de objetos em atividades fotocopiadas com desenhos prontos ou de quantidades em \textbf{que} elas devem ligar o número ao conjunto de objetos correspondentes.
Essas atividades não possibilitam a compreensão efetiva dessas noções pela criança e nem condizem com o conceito de criança apresentado nos documentos legais, como as DCNEI ( BRASIL, 2009 ) e a BNCC ( BRASIL, 2017 ).
Os conhecimentos matemáticos abrangem : espaço e forma ; números ; grandezas e medidas e estatística. \textbf{Contudo}, essa é \textbf{uma} \textbf{divisão} didática para o(a) professor(a), \textbf{que}, no trabalho cotidiano, deve propor situações de aprendizagens, articulando vários conhecimentos, inclusive os de outros campos de experiências.
O desenvolvimento de noções relacionadas às formas, planas e tridimensionais, envolve a \textbf{percepção} e a compreensão de suas características - tamanho, forma, cor, textura - bem como, a sua localização no espaço no \textbf{que} se refere à \textbf{direção} e posição.
O aprendizado desses conhecimentos se dá em situações de : exploração e manipulação de objetos ; brincadeiras com blocos lógicos ou de construção ; modelagem com massinhas variadas ; leitura de histórias ; organização do espaço da sala ; produção de figuras a partir do tangram ; brincadeiras em espaços amplos ; passeios pelo entorno para observação de figuras presentes em fachadas e construções ; apreciação de obras de arte etc.
Nesses contextos, de acordo com Smole, Diniz e Cândido ( 2003 ), o professor deve fazer mediações com perguntas, \textbf{problematizações} e uso de vocabulários e conceitos adequados, para \textbf{que} a criança possa : coordenar sua visão com o movimento do corpo, como a brincadeira de amarelinha ; \textbf{lembrar} de \textbf{um} objeto \textbf{que} não está no seu campo de visão e distinguir suas semelhanças e diferenças ao serem relacionados com outros ; identificar figuras planas em diferentes contextos - edificações, objetos etc. ; relacionar dois objetos entre si - se estão próximos ou distantes, qual é maior ou menor, o \textbf{que} está em cima ou embaixo etc. ; conservar propriedades - independente da localização do objeto o seu tamanho não é alterado.
O conceito de número é \textbf{bastante} complexo.
A criança, na tentativa de entender o seu significado e a maneira de utilizá -los no dia a dia, elabora suas próprias \textbf{teorias} e explicações, \textbf{que} são modificadas ao longo do tempo, à medida \textbf{que} é desafiada a pensar e ampliar seus conhecimentos sobre eles.
Por isso, a crítica quanto a afirmação de \textbf{que} a compreensão de conceitos matemáticos ocorre somente pela manipulação de material concreto, \textbf{uma} vez \textbf{que} a ação do \textbf{sujeito} envolve e pressupõe o pensamento.
Participar de situações de aprendizagens em \textbf{que} são faladas ou cantadas parlendas, músicas e brincadeiras \textbf{que} tem a recitação ordenada de números, auxilia a criança a memorizar a série numérica \textbf{que}, posteriormente, favorecerá a compreensão de regras \textbf{que} regem o Sistema de Numeração Decimal e do \textbf{que} é quantificar – contar \textbf{um} conjunto de elementos sem repetir ou saltar nenhum e entender \textbf{que} o último número representa todo o conjunto e não só o último elemento.
Recitar não significa a mesma coisa \textbf{que} contar ou quantificar.
Crianças \textbf{dizem} a série ordenada de números sem necessariamente compreender o \textbf{que} é quantidade, porque a recitação não pressupõe \textbf{uma} situação de enumeração.
Já na contagem é \textbf{preciso} estabelecer a correspondência termo a termo ou \textbf{um} a \textbf{um}, \textbf{que} é relacionar o nome do número ao elemento a ser contado, para \textbf{que} se tenha \textbf{uma} quantidade.
As contagens devem fazer parte do dia a dia da instituição de Educação Infantil, como : na identificação de quantos meninos e meninas estão presentes ; na distribuição de alimentos pela cozinheira para \textbf{uma} turma ou agrupamento ; na organização de materiais para realização de brincadeiras ; no acompanhamento das coleções de figurinhas, de carrinhos, entre outras.
Nesses contextos, é \textbf{preciso} o professor incentivar a criança a fazer os registros de forma não convencional com desenhos, pauzinhos, bolinhas e convencional com números, assim como, dialogar com os colegas sobre as estratégias utilizadas para chegar ao resultado.
O quadro numérico²² pode ser \textbf{um} importante aliado da criança na consulta da escrita dos números.
Perguntas como : “ Quem tem mais? ”, “ Quem tem menos? ”, “ Quantos têm o mesmo tanto? ”, “ Quem tem quantidades parecidas? ”, “ Quem tem quantidades diferentes? ”, “ Quantos são iguais? ”, na resolução de problemas e na contagem de diferentes elementos também favorecem a criança a pensar sobre quantidades.
O trabalho com números, em contextos significativos, pressupõe abordar, também, grandezas e medidas, relacionadas à : distância – comprimento, espessura, altura, largura, tamanho, profundidade ; superfície – brilho, cor, textura etc. ; espaço – capacidade, volume ; massa – matéria ou peso como é \textbf{conhecida} ; calor ; movimento – velocidade e ; duração – regularidades, ritmos ou alternâncias ( LORENZATO, 2006 ).
Para cada grandeza existe \textbf{uma} forma de medir e instrumentos próprios \textbf{que} permitem essa quantificação, como : trenas, fitas métricas, copos medidores, relógios, termômetros, velocímetros, calendários etc.
A presença e o uso desses materiais pela criança no dia a dia são importantes para : comparar alturas e pesos ; apropriar das regularidades do tempo ( dia, mês e ano ) ; aprender a fazer leitura em relógio digital e analógico ; produzir receitas ; medir a sala para reorganizar o mobiliário ou \textbf{uma} área externa para fazer canteiros de \textbf{uma} horta etc.
O sistema monetário, também, deve ser trabalhado em situações lúdicas e de faz de conta, para \textbf{que} a criança possa identificar e manusear cédulas e moedas, fazendo associações com preços.
A estatística, por utilizar \textbf{uma} linguagem visual, possibilita aos \textbf{sujeitos} lerem de forma mais rápida sínteses e conclusões referentes a \textbf{um} determinado assunto por meio de tabelas e gráficos.
As ações envolvidas nesse processo são a coleta, a organização e a síntese de informações com intuito de fazer comparações e análises \textbf{que} favoreçam o desenvolvimento da capacidade argumentativa e do pensamento analítico.
As tabelas e gráficos podem ser elaborados e explorados em suas várias formas e com assuntos diversos, como : se a instituição tem mais meninos ou meninas ; os números dos sapatos ; as medidas do corpo ; as preferências das crianças no \textbf{que} se refere a alimentos, a brinquedos, a brincadeiras, a jogos ; em pesquisas de opinião sobre a escolha do nome de \textbf{um} projeto ou em \textbf{que} a verba da instituição pode ser investida etc.
Essas ações devem ser articuladas aos projetos de trabalho, temáticos ou investigativos de forma contextualizada, previstas no plano diário dos(as) professores a partir de situações lúdicas e criativas com a participação e o envolvimento efetivo das crianças para darem sentido ao \textbf{que} está sendo realizado.

% matched lemmas: ator, basear, bastante, central, conhecido, conseguir, contemporaneidade, contudo, correlação, depender, desempenhar, direção, divisão, dizer, hoje, homem, humanidade, imediato, implicar, influência, inúmero, lado, lembrar, marca, operar, percepção, permear, pontuar, preciso, problematização, que, remeter, ribeirinho, sistematizar, sujeito, tarefa, teoria, um, urgente, ver, viver
\end{textsample}
